%! Graphing Quadratic Functions & Transformations — Unit 3, Lesson 3.1 — Group Activity
\documentclass[10pt]{article}
\usepackage{algebra2-article}
\usepackage{algebra2-boxes}
\newcommand{\TallMath}[1]{$\displaystyle #1\rule[-1.4em]{0pt}{3.2em}$}
\newcommand{\lgrid}[4]{%
  \draw[step=1, linegray!40, very thin] (#1,#3) grid (#2,#4);
  \draw[->, semithick, charcoal] (#1,0)--(#2,0) node[below right, font=\scriptsize]{$x$};
  \draw[->, semithick, charcoal] (0,#3)--(0,#4) node[left, font=\scriptsize]{$y$};}

\begin{document}

\pageheader{Unit 3, Lesson 3.1}{Group Activity}
\namepartnerperiod

\vspace{0.08in}

\begin{headlinebox}{forestbg}
\small\textbf{Three Costumes, One Curve.} A quadratic can be written in vertex, standard, or intercept
form --- each one showing off a different feature. With your partner, pull the features out of each
form, then use them to place or graph the parabola. \emph{Justify} each vertex and axis.
\end{headlinebox}

\vspace{0.06in}

\begin{tcolorbox}[colback=white, colframe=black!40, title=\textbf{Tier R --- Remediate},
    fonttitle=\bfseries, arc=1.5mm, left=3mm, right=3mm, top=2mm, bottom=2mm, breakable]
Each equation is in \textbf{vertex form} $a(x-h)^2+k$. Read the vertex $(h,k)$ and the opening direction,
then match it to graph A, B, or C.
\begin{center}
\begin{tikzpicture}[scale=0.38]
  % A: (x-3)^2 - 2  vertex (3,-2) up
  \begin{scope}
    \lgrid{0}{6}{-3}{3}
    \begin{scope}\clip (0,-3) rectangle (6,3);
      \draw[very thick, forest, domain=1.3:4.7, samples=60, smooth] plot (\x,{(\x-3)*(\x-3) - 2});
    \end{scope}
    \fill[charcoal] (3,-2) circle (2.5pt);
    \node[charcoal, font=\scriptsize] at (3,-4.2){\textbf{A}};
  \end{scope}
  % B: -(x+1)^2 + 2  vertex (-1,2) down
  \begin{scope}[xshift=8cm]
    \lgrid{-4}{2}{-3}{3}
    \begin{scope}\clip (-4,-3) rectangle (2,3);
      \draw[very thick, forest, domain=-3.2:1.2, samples=60, smooth] plot (\x,{-1*(\x+1)*(\x+1) + 2});
    \end{scope}
    \fill[charcoal] (-1,2) circle (2.5pt);
    \node[charcoal, font=\scriptsize] at (-1,-4.2){\textbf{B}};
  \end{scope}
  % C: (x+2)^2 - 1  vertex (-2,-1) up
  \begin{scope}[xshift=16cm]
    \lgrid{-5}{1}{-2}{4}
    \begin{scope}\clip (-5,-2) rectangle (1,4);
      \draw[very thick, forest, domain=-4.2:-0.3, samples=60, smooth] plot (\x,{(\x+2)*(\x+2) - 1});
    \end{scope}
    \fill[charcoal] (-2,-1) circle (2.5pt);
    \node[charcoal, font=\scriptsize] at (-2,-3.2){\textbf{C}};
  \end{scope}
\end{tikzpicture}
\end{center}
\renewcommand{\arraystretch}{1.5}
\begin{tabularx}{\linewidth}{@{}l X X X c@{}}
\textbf{Equation} & \textbf{Vertex} & \textbf{Axis} & \textbf{Opens} & \textbf{Graph} \\
\hline
$y=(x-3)^2-2$   & \blank{1.8cm} & $x=\blank{0.9cm}$ & \blank{1.2cm} & \blank{0.7cm} \\
$y=-(x+1)^2+2$  & \blank{1.8cm} & $x=\blank{0.9cm}$ & \blank{1.2cm} & \blank{0.7cm} \\
$y=(x+2)^2-1$   & \blank{1.8cm} & $x=\blank{0.9cm}$ & \blank{1.2cm} & \blank{0.7cm} \\
\end{tabularx}
\end{tcolorbox}

\begin{tcolorbox}[colback=white, colframe=black!40, title=\textbf{Tier A --- Approaching Proficiency},
    fonttitle=\bfseries, arc=1.5mm, left=3mm, right=3mm, top=2mm, bottom=2mm, breakable]
\textbf{1. Standard form.} For $f(x)=x^2-2x-3$ ($a=1,\ b=-2,\ c=-3$):
\begin{enumerate}[leftmargin=1.7em, itemsep=4pt]
  \item Axis of symmetry $x=-\dfrac{b}{2a}=\blank{1.2cm}$; \; vertex $\blank{1.8cm}$; \; $y$-intercept $\blank{1.6cm}$.
  \item Complete the table (notice it is symmetric about the axis).
  \begin{center}
  \renewcommand{\arraystretch}{1.3}
  \begin{tabular}{|c|c|c|c|c|c|}
  \hline
  $x$ & $-1$ & $0$ & $1$ & $2$ & $3$ \\ \hline
  $f(x)$ & \blank{0.8cm} & \blank{0.8cm} & \blank{0.8cm} & \blank{0.8cm} & \blank{0.8cm} \\ \hline
  \end{tabular}
  \end{center}
\end{enumerate}
\textbf{2. Convert vertex $\to$ standard.} Expand $y=(x-2)^2-5$ into the form $x^2+bx+c$:
\;\blank{4.5cm}
\end{tcolorbox}

\begin{tcolorbox}[colback=white, colframe=black!40, title=\textbf{Tier E --- Extension},
    fonttitle=\bfseries, arc=1.5mm, left=3mm, right=3mm, top=2mm, bottom=2mm, breakable]
\textbf{1. Intercept form.} For $h(x)=(x-1)(x-5)$: zeros $x=\blank{0.9cm}$ and $x=\blank{0.9cm}$; axis
$x=\blank{0.9cm}$; vertex $\blank{1.8cm}$; $y$-intercept $\blank{1.6cm}$.

\vspace{4pt}
\textbf{2. Write the equation from the graph.} The parabola below has vertex $(-2,-4)$ and passes
through the origin.
\begin{center}
\begin{tikzpicture}[scale=0.45]
  \lgrid{-6}{2}{-5}{3}
  \draw[navy, dashed, semithick] (-2,-5)--(-2,3);
  \begin{scope}\clip (-6,-5) rectangle (2,3);
    \draw[very thick, forest, domain=-4.4:0.4, samples=70, smooth] plot (\x,{(\x+2)*(\x+2) - 4});
  \end{scope}
  \fill[charcoal] (-2,-4) circle (3pt) node[below right, font=\scriptsize]{$(-2,-4)$};
  \fill[charcoal] (-4,0) circle (3pt);
  \fill[charcoal] (0,0) circle (3pt);
\end{tikzpicture}
\end{center}
Using vertex form $y=a(x-h)^2+k$ with vertex $(-2,-4)$ and the point $(0,0)$, solve for $a$ and write the
equation. \writelines{2}

\vspace{2pt}
\textbf{3. Describe the moves.} List the transformations, \emph{in order}, that carry $y=x^2$ onto
$y=-2(x-1)^2+3$. \writelines{2}
\end{tcolorbox}

\end{document}
